\documentclass[11pt]{article}
\usepackage[margin=1in]{geometry}
\usepackage[T1]{fontenc}
\usepackage[utf8]{inputenc}
\usepackage{lmodern}
\usepackage{hyperref}
\usepackage{graphicx}
\usepackage{booktabs}
\usepackage{enumitem}
\usepackage{xcolor}
\usepackage{longtable}
\hypersetup{
  colorlinks=true,
  linkcolor=blue,
  urlcolor=blue,
  pdftitle={Foreign Whispers Final Report},
  pdfauthor={Ahmed Jaheen}
}

\title{Foreign Whispers:\\Final Project Report}
\author{Ahmed Jaheen}
\date{\today}

\begin{document}
\maketitle

\begin{abstract}
This report presents the implementation of \textit{Foreign Whispers}, an
open-source video dubbing pipeline that transforms an English YouTube video
into a dubbed Spanish output using local components. The final system supports
video download, transcription, translation, speaker diarization, speaker-aware
voice selection, time-aware text-to-speech, and final audio/video stitching.
The project was completed with an emphasis on practical usability: the codebase
now runs through a browser-based studio, exposes a FastAPI backend, supports a
Docker GPU workflow, and also includes a local macOS CPU fallback path for
development and validation. This report summarizes the architecture, the main
technical decisions, the implemented notebook tasks, debugging work, and the
final deliverables.
\end{abstract}

\section{Project Goal}
The goal of Foreign Whispers is to recreate the essential stages of a modern
video dubbing system using open-source tools and local execution instead of
paid APIs. The pipeline accepts a YouTube URL, downloads the source video and
captions, transcribes spoken audio, translates the transcript into Spanish,
synthesizes a dubbed audio track, and stitches that audio back into the source
video while generating subtitle tracks for playback in the frontend.

Compared with a simple translation demo, this project is harder because each
stage has timing consequences for the next stage. A translation that is
semantically correct can still fail if it is too long to be spoken in the
original window. Likewise, text-to-speech audio that sounds natural can still
produce a poor video if it drifts away from the visible speaker. The final
implementation therefore focuses not only on correctness at each stage, but
also on timing stability across the entire pipeline.

\section{System Architecture}
The final system uses a layered architecture with a browser frontend, a FastAPI
orchestrator, optional GPU inference services, and a reusable Python library.

\subsection{Major Components}
\begin{itemize}[leftmargin=*]
\item \textbf{Frontend (Next.js):} The Dubbing Studio runs on port 8501 and
provides video selection, pipeline settings, stage monitoring, result playback,
and caption rendering.
\item \textbf{Backend (FastAPI):} The API runs on port 8080 and exposes the
pipeline routes for download, transcription, diarization, translation, TTS,
stitching, and media serving.
\item \textbf{Library (\texttt{foreign\_whispers}):} This package contains the
alignment logic, evaluation helpers, reranking code, diarization helpers, and
voice resolution functions. It is intentionally reusable outside the web app.
\item \textbf{Inference Services:} In the intended GPU deployment, Whisper and
Chatterbox run in dedicated services. In local macOS development, the backend
can fall back to CPU-safe implementations.
\item \textbf{Artifact Cache:} All pipeline outputs are stored under
\texttt{pipeline\_data/api/}. This makes repeated runs faster and supports
inspection of each intermediate stage.
\end{itemize}

\subsection{Pipeline Flow}
The full processing sequence is:
\begin{enumerate}[leftmargin=*]
\item Download the YouTube video and captions.
\item Transcribe the video audio using Whisper.
\item Diarize the audio to determine which speaker is talking in each region.
\item Translate the transcript from English to Spanish.
\item Resolve voices and synthesize Spanish audio with timing control.
\item Stitch the dubbed audio and WebVTT subtitles back into the final video.
\end{enumerate}

The frontend calls the backend step by step, and the backend stores the
artifacts of every stage so that debugging and reruns are manageable.

\section{Implementation Details}
\subsection{Download and Transcription}
The download stage uses \texttt{yt-dlp} to retrieve the source MP4 and YouTube
captions. The transcription stage uses Whisper. One important practical issue
was that YouTube captions are often too coarse for dubbing because they cover
long windows rather than short spoken phrases. For that reason, the frontend
defaults to Whisper timing instead of YouTube caption timing.

Additional segmentation logic was used to split overlong transcription windows
into smaller phrase-like segments. This significantly improved alignment and
reduced the silence-heavy behavior observed in early runs.

\subsection{Translation and Duration-Aware Reranking}
The translation stage uses Argos Translate for offline English-to-Spanish
translation. However, direct translation alone is not enough, because Spanish
often expands relative to English. To handle this, the project implements
duration-aware translation reranking in \texttt{foreign\_whispers/reranking.py}.

The reranking helper generates shorter alternatives when a translated segment
is too long for its target timing window. This improvement is reused in both
the translation notebook and the alignment notebook. In practice, it reduced
the number of severe stretch cases and helped prevent missing words in TTS.

\subsection{Alignment and Timing Control}
The alignment work was one of the most important parts of the project. The
initial heuristic estimated TTS duration mostly from character count, which was
too crude. The improved implementation uses a better duration estimate and a
stronger global alignment strategy to schedule translated segments more
intelligently.

Another important refinement was distinguishing \textit{natural pacing} from
\textit{forced stretching}. In baseline mode, the system now preserves natural
speech and allows drift. In aligned mode, it tries to fit speech to the video
window while avoiding overly slow or heavily clipped audio. Short segments are
no longer artificially slowed just to fill the whole subtitle window; instead,
the system prefers silence padding when that sounds better.

\subsection{Diarization and Speaker Assignment}
The diarization stage was implemented with \texttt{pyannote.audio}. This stage
extracts audio from the source video, runs diarization, caches the speaker
segments, and merges speaker labels back into the transcript. The core merging
logic assigns each transcription segment to the speaker with the strongest
temporal overlap.

Because the project also needed male and female speaker-aware dubbing, the
implementation builds lightweight speaker profiles from pitch estimates over
the diarized audio regions. Those profiles are not full biometric models; they
are simple cues used to resolve an appropriate male or female reference voice.
This is sufficient for the course requirement that a woman should sound like a
woman and a man should sound like a man.

\subsection{Voice Resolution and TTS}
Voice resolution was implemented in
\texttt{foreign\_whispers/voice\_resolution.py}. The resolution order is:
\begin{enumerate}[leftmargin=*]
\item speaker-specific WAV
\item language-specific gender default
\item language default
\item global default
\end{enumerate}

This design makes the system robust. When a dedicated speaker reference file is
available, it is used. Otherwise, the pipeline can still choose a male or
female language default. For Spanish, reference WAV files were added under
\texttt{pipeline\_data/speakers/es/}.

In the intended Docker GPU setup, TTS is served by Chatterbox. In local macOS
development, however, Chatterbox may not be available. To keep the project
fully runnable, CPU-safe fallbacks were added. A major late-stage bug was that
the local fallback used a single generic voice, which made correct diarization
sound wrong to the user. This was fixed by routing male and female diarized
segments to distinct local macOS Spanish voices when Chatterbox is unavailable.

\subsection{Stitching and Captions}
The stitch stage remuxes the synthesized audio back into the source MP4 using
\texttt{ffmpeg} without re-encoding the video track. This keeps the output
fast and high quality. The system also serves WebVTT subtitle files for both
original and translated captions.

A caption rendering bug was discovered where translated text appeared twice on
top of itself in the player. The root cause was overlapping cue behavior from a
carry-over caption format. The stitch code was updated so the frontend now
receives clean one-cue subtitle timing for the translated captions.

\section{Debugging and Engineering Improvements}
Several practical issues had to be solved before the application became stable.

\subsection{Cache Invalidation}
The pipeline caches intermediate results aggressively. This is useful for
performance, but it created a class of bugs where later code improvements did
not appear in the output because old translation or TTS files were reused. The
backend was updated to invalidate caches when speaker labels or speaker voice
metadata were missing from downstream artifacts.

\subsection{Frontend Pipeline State}
Another issue appeared in the frontend stage table: runs could look as if
multiple stages were executing out of order. The root cause was overlapping or
stale asynchronous updates from different runs. The pipeline hook was fixed so
only the current run can update stage state, and overlapping runs are now
blocked.

\subsection{Settings Dependency Handling}
The user also observed that the diarization stage was being skipped even after
selecting speaker-aware TTS. This came from a frontend settings dependency
problem. The fix was to make diarization a derived requirement when voice
cloning is enabled, so the UI and pipeline logic stay consistent.

\section{Evaluation and Testing}
The project was validated with both automated tests and real pipeline runs.
The backend test suite passes successfully, and frontend lint/build checks were
also used after UI changes. In addition, several end-to-end reruns were used to
verify that early timeline regions no longer contained long silence windows,
that subtitle duplication was fixed, and that speaker-aware male/female voice
assignment worked in practice.

The final codebase therefore supports both analytical testing and practical
user-facing verification through the browser interface.

\section{Deliverables}
\subsection{Codebase and README}
The repository now contains:
\begin{itemize}[leftmargin=*]
\item the full application code
\item exact Docker and macOS local run instructions
\item environment variable guidance
\item output artifact locations
\item validation commands for tests, lint, and build checks
\end{itemize}

\subsection{Report}
This LaTeX file is intended as the report deliverable. It can be copied
directly into Overleaf and expanded with screenshots or experiment tables if
desired.

\subsection{Demo Video and Sample Input/Output}
The demo video and sample input/output videos are separate deliverables. The
repo already contains generated outputs under \texttt{pipeline\_data/api/},
which can be used as references during recording.

\section{Conclusion}
The final Foreign Whispers implementation goes beyond a basic proof of concept.
It now behaves like a usable end-to-end dubbing system with a frontend studio,
an API backend, reusable alignment logic, diarization support, speaker-aware
voice selection, caching, and stable local development workflows.

The most important lesson from the project is that dubbing quality depends on
the interaction between components, not just the quality of any one model.
Translation length, segment timing, speaker assignment, and TTS behavior all
affect the final video. The completed system addresses these interactions
through alignment-aware translation, cache correctness, speaker-aware fallback
voices, and a more reliable frontend execution model.

\end{document}
